\begin{center}\begin{tikzpicture}[font=\small,>=stealth,line join=round]
\draw[gray!40,rounded corners=3pt] (0,0) rectangle (14.2,8.3);\node at (7.1,7.700000000000001) {证明三：对数在1的两侧异号，决定比较方向};
\begin{scope}[shift={(1.7000000000000002,4.205555555555556)},x=1.5cm,y=1.0222222222222221cm]\draw[->,gray] (-0.4,0)--(3.6,0) node[right]{$x$};\draw[->,gray] (0,-2.5)--(0,2) node[above]{$y$};\draw[blue,thick,solid,domain=0.12:3.2,samples=180] plot(\x,{ln(\x)});\fill[orange!85!black] (0.5,-0.6931471805599453) circle(2.2pt);\fill[orange!85!black] (1,0) circle(2.2pt);\fill[orange!85!black] (2,0.6931471805599453) circle(2.2pt);\draw[gray,dashed] (1,0)--(1,1.6);\node[below right=6pt] at (.5,-.693147){$\ln\frac12<0$};\node[below right=4pt] at (1,0){$1$};\node[above=10pt] at (2,.693147){$\ln2>0$};\end{scope}\node[anchor=west,align=left,text width=5.0cm] at (8.4,4.0) {$\ln(x^{\alpha_1})-\ln(x^{\alpha_2})$\\[9pt]$=(\alpha_1-\alpha_2)\ln x$\\[9pt]若$\alpha_1>\alpha_2$：\\[9pt]差的符号由$\ln x$决定\\[9pt]过1时：负、零、正};
\node[text width=12.6cm,align=center] at (7.1,.65) {蓝线是ln x；对数严格递增，所以对数差的符号也就是原来两数的大小方向。};
\end{tikzpicture}\end{center}
